\documentclass[11pt]{article}

% Required Packages
\usepackage[utf8]{inputenc}
\usepackage{graphicx}       % For handling images
\usepackage{amsmath}        % For advanced math (cases, align)
\usepackage{amssymb}        % For math symbols
\usepackage{subcaption}     % For sub-figures (Figure 1a, 1b, 1c)
\usepackage{multirow}       % For table formatting if needed

% Page setup (Optional: Adjust margins to look like a standard assignment)
\usepackage[margin=1in]{geometry}

\title{Mathematical Foundations of Image Synthesis}
\author{Your Name (Your Student ID)}
\date{January 4, 2026}

\begin{document}

\maketitle

\section{Introduction}
Image synthesis is the process of generating images from mathematical and algorithmic descriptions. Modern graphics systems rely on precise mathematical models and efficient computation. Even small numerical errors can result in visible artifacts in rendered images. Typical formulations involve scalar values such as $\alpha$, indexed parameters like $I_{ij}$, and exponential terms $e^{x^{2}}$. These elements are combined to simulate realistic lighting behavior.

\section{Core Components}

\subsection{Rendering Pipeline}
The rendering pipeline consists of the following stages:
\begin{itemize}
    \item Scene Definition
    \item Camera Configuration
    \item Light Placement
    \item Important Shading Model
    \begin{itemize}
        \item Diffuse Shading
        \item Specular Highlights
    \end{itemize}
\end{itemize}

\subsection{Execution Order}
\begin{enumerate}
    \item Input Processing
    \item Ray Construction
    \begin{enumerate}
        \item[(a)] Primary Rays
        \item[(b)] Secondary Rays
    \end{enumerate}
    \item Intersection Evaluation
    \begin{itemize}
        \item Object Hit Test
    \end{itemize}
\end{enumerate}

\clearpage  % Sections 3 and 4 start here on their own page

\section{Mathematical Modeling}
Light interaction with surfaces can be expressed using the following equations.

\begin{equation}
    L = I_{ij} \cdot \cos(\theta) + R^{2}
\end{equation}

\begin{equation}
    k = \frac{ab}{cd}
\end{equation}

\begin{equation}
\begin{split}
    f(x) & = x^{2} + 2x + 1 \\
         & = (x+1)^{2}
\end{split}
\end{equation}

\begin{equation}
    |y| = \begin{cases} 
        y & \text{if } y \ge 0 \\
        -y & \text{if } y < 0 
    \end{cases}
\end{equation}

\begin{equation}
    p(t) = \begin{cases} 
        0, & t < 0, \\
        e^{-t}, & 0 \le t < 3, \\
        \ln(t-2), & t \ge 3
    \end{cases}
\end{equation}

\begin{equation}
    \mathcal{N}(x_{i},\mu,\sigma) = \frac{1}{\sqrt{2\pi}\sigma}e^{-\frac{(x_{i}-\mu)^{2}}{2\sigma^{2}}}
\end{equation}

\section{Performance Comparison}

The following table:

\begin{table}[h!]
    \centering
    \caption{Comparison of Image Synthesis Techniques}
    \label{tab:comparison}
    \vspace{0.2cm} % Add a little space between caption and table
    \begin{tabular}{|l|l|l|}
        \hline
        \textbf{Method} & \textbf{Accuracy} & \textbf{Time Complexity} \\ \hline
        Rasterization & Medium & $O(n)$ \\ \hline
        Ray Tracing & \begin{tabular}{@{}l@{}}High \\ Very High\end{tabular} & \begin{tabular}{@{}l@{}}$O(n^{2})$ \\ $O(n^{3})$\end{tabular} \\ \hline
        Hybrid Method & & $O(n^{2})$ \\ \hline
    \end{tabular}
\end{table}

The results in Table \ref{tab:comparison} highlight the trade-off between accuracy and computational cost.

\clearpage  % Sections 5 and 6 start here on the next page

\section{Visual Illustration}

Figure \ref{fig:rendering_outputs} compares three different rendering outputs arranged in a single row.

\begin{figure}[h!]
    \centering
    \begin{subfigure}[b]{0.3\textwidth}
        \centering
        \includegraphics[width=\textwidth]{example-image-a}
        \caption{Output A}
    \end{subfigure}
    \hfill
    \begin{subfigure}[b]{0.3\textwidth}
        \centering
        \includegraphics[width=\textwidth]{example-image-b}
        \caption{Output B}
    \end{subfigure}
    \hfill
    \begin{subfigure}[b]{0.3\textwidth}
        \centering
        \includegraphics[width=\textwidth]{example-image-c}
        \caption{Output C}
    \end{subfigure}
    \caption{Comparison of three rendering outputs.}
    \label{fig:rendering_outputs}
\end{figure}

\section{Conclusion}
Image synthesis combines mathematical precision, algorithmic design, and efficient implementation. As scene complexity increases, choosing the appropriate method becomes critical for balancing quality and performance.

\end{document}